\section{A compact current-reference architecture}
\label{sec:controller}

The recommended linear-model implementation combines offline/pre-update matrix
construction with a scalar online boundary search.

\subsection{Preparation at a parameter update}

Given
\(\theta=(R_s,R_{\mathrm{exc}},L_d,L_q,L_{\mathrm{exc}},p)\) and speed,
form \(\matr R,\matr H,\matr Q_U\), then
\(\bar{\matr H}=\matr R^{-1/2}\matr H\matr R^{-1/2}\) and
\(\bar{\matr Q}=\matr R^{-1/2}\matr Q_U\matr R^{-1/2}\). Temperature-updated
resistances may enter here at runtime. Cache the six independent entries of
each symmetric matrix and the inverse whitening scale.

\subsection{Online request handling}

\begin{enumerate}
 \item Select the torque-sign branch. Evaluate the unconstrained best-loss
       direction from the extreme eigenvector of \(\bar{\matr H}\).
 \item Scale it to \(M_{\mathrm{ref}}\) and test voltage. If voltage is
       inactive, compute \(P_{\mathrm{req}}\) directly.
 \item If voltage is active, bracket the support slope \(\mu\). At each trial,
       compute the dominant eigenvector of
       \(\bar{\matr H}-\mu\bar{\matr Q}\), map it to \((\nu,\tau)\), and use
       the signed residual
       \(r(\mu)=M_{\mathrm{ref}}\nu-U_{\max}^2\tau\).
       Bisection or safeguarded interpolation finds the boundary contact.
 \item Set \(P_{\mathrm{req}}=M_{\mathrm{ref}}/\tau_\star\). If
       \(P_{\mathrm{req}}\leq P_{\mathrm{Cu},\max}\), scale
       \(\vect y_\star=\sqrt{P_{\mathrm{req}}}\vect x_\star\), back-transform,
       and return status \texttt{TRACKING}.
 \item Otherwise maximize \cref{eq:performance-mmax} on the same cached upper
       boundary, scale to the first active resource, and return status
       \texttt{LOSS\_LIMITED}, \texttt{VOLTAGE\_LIMITED}, or
       \texttt{BOTH\_LIMITED}. Degenerate parameters or empty positive-torque
       branches return \texttt{MODEL\_INVALID} or \texttt{INFEASIBLE}.
\end{enumerate}

The interface is
\begin{equation}
 (i_d^\star,i_q^\star,i_{\mathrm{exc}}^\star,M_{\mathrm{ach}},
 P_{\mathrm{Cu}}^\star,\mathrm{status})
 =\gls{sym:currentmap}(M_{\mathrm{ref}},U_{\max},P_{\mathrm{Cu},\max},
 \omega_e;\theta).
 \label{eq:controller-map}
\end{equation}

\subsection{Cost and numerical choices}

A 3D symmetric matrix requires six stored scalars. Two whitened forms, the
back-transform, brackets, and status data fit comfortably below 32 floating
values. The voltage-inactive branch needs one precomputed eigenvector, one
division, one square root, and a matrix-vector back-transform online. The
voltage-active branch typically uses 8--12 bracket iterations for a fixed
relative tolerance. Each trial forms six affine matrix entries and solves one
symmetric \(3\times3\) extreme eigenpair. A compact Jacobi implementation
using 5--8 plane rotations is roughly \(80\)--\(150\) multiplications per
trial, with square roots or hypot operations in the rotations; exact counts
depend on whether divisions and rotations are tabled, CORDIC-based, or
hardware-supported. No trigonometric or hyperbolic function is required.

For an \gls{mcu}, fixed-iteration bisection plus a robust symmetric eigensolver
is easier to certify than radicals for a sextic contact equation. For an
\gls{fpga}, the three Jacobi rotation cells can be pipelined; alternatively a
precomputed support-slope grid plus one interpolation supplies deterministic
latency. Coefficient scaling should normalize \(\bar{\matr H}\) and
\(\bar{\matr Q}\) near unity. Eigenvector signs are chosen continuously using
the previous cycle, then corrected to the requested torque branch. Every
returned point is checked in the original equations, not only in transformed
coordinates.

This is a structured numerical method rather than a fetish for symbolic
closed form. Its iterations are one-dimensional, monotone/bracketed, and tied
to visible geometry; their bounded cost is usually preferable to fragile
general quartic or sextic root selection.
